\section{Requisiti Funzionali}
\label{sec:RF}

In questa sezione sono elencati i requisiti funzionali di TEPO, organizzati per area funzionale. Ogni requisito descrive una capacità o un comportamento richiesto al sistema, le informazioni coinvolte e le principali regole di dominio note. La sequenza delle interazioni utente--sistema è dettagliata nella Sezione~\ref{sec:UC} (Use Case); le proprietà qualitative del sistema sono invece trattate nella Sezione~\ref{sec:RNF} (Requisiti Non Funzionali). Gli identificatori \texttt{RF-01}, \texttt{RF-02}, \ldots, \texttt{RF-19} sono stabili a partire da questa baseline e vengono mantenuti nella matrice di tracciabilità e negli artefatti successivi.

\subsection{Requisiti funzionali di catalogo e consultazione}

\RF{Consultazione del catalogo degli edifici.} \label{rf:catalogo_edifici}
Il sistema deve consentire all'utente di consultare il catalogo degli edifici pubblici di interesse, mostrando per ciascun edificio: identificativo, denominazione, indirizzo, comune, tipologia d'uso, superficie, numero di interventi candidati associati e provenienza del dato (reale, sintetica o mista). Il catalogo deve poter essere filtrato per tipologia d'uso e ricercato per nome, indirizzo, comune o identificativo; le colonne principali devono poter essere ordinate in modo crescente e decrescente. Il sistema deve inoltre mostrare indicatori sintetici del catalogo (numero di edifici, numero di interventi collegati, superficie totale e tipologie rappresentate) e deve rendere esplicito quando un filtro non produce risultati.

\RF{Visualizzazione geografica degli edifici.} \label{rf:mappa}
Il sistema deve offrire una vista cartografica degli edifici del catalogo, posizionando ciascun edificio alle coordinate registrate. La vista deve permettere di distinguere lo stato di ogni edificio rispetto all'ultimo risultato di ottimizzazione disponibile: con interventi selezionati, con interventi selezionati ed esclusi, con soli interventi esclusi, non valutato. Lo stato non deve essere veicolato dal solo colore: il sistema deve affiancare al colore una forma o un contorno del marcatore e un'etichetta testuale, come stabilito da \RNFref{rnf:accessibilita}. La dimensione del marcatore deve rappresentare l'investimento associato all'edificio. La vista deve offrire una legenda, la ricerca degli edifici e un elenco accessibile equivalente alla mappa; selezionando un edificio dall'elenco o dalla mappa, il sistema deve mostrarne il dettaglio (indirizzo, comune, tipologia, superficie, stato e interventi associati).

\RF{Consultazione del catalogo degli interventi.} \label{rf:catalogo_interventi}
Il sistema deve consentire di consultare l'insieme delle opzioni di intervento disponibili, mostrando per ciascuna: edificio di riferimento, denominazione e descrizione, tipologia di intervento (illuminazione LED, serramenti, isolamento, pompa di calore, fotovoltaico, caldaia), costo stimato, risparmio energetico annuo atteso, risparmio economico annuo atteso, riduzione annua attesa di CO\textsubscript{2}, payback semplice e provenienza del dato. Il payback deve essere etichettato come stima demo e non come misura ufficiale. Il sistema deve inoltre mostrare lo stato di ciascuna opzione rispetto all'ultimo risultato eseguito (selezionato, escluso, non valutato) e gli indicatori sintetici del sottoinsieme filtrato.

\RF{Ricerca, filtraggio e ordinamento degli interventi.} \label{rf:filtri_interventi}
Il sistema deve permettere di ricercare e filtrare le opzioni di intervento combinando più criteri: ricerca testuale, edificio di riferimento, tipologia di intervento, intervallo di costo e stato rispetto all'ultimo risultato. I filtri devono essere combinabili e devono poter essere azzerati; il sistema deve mostrare in ogni momento il numero di interventi che soddisfano i criteri applicati. L'utente deve poter ordinare i risultati per le principali grandezze numeriche e testuali (identificativo, edificio, costo, risparmio energetico, risparmio economico, CO\textsubscript{2}, payback) in entrambe le direzioni. Se nessun intervento soddisfa i criteri, il sistema deve mostrare un messaggio informativo con la possibilità di azzerare la ricerca.

\RF{Importazione di dati di catalogo da file CSV.} \label{rf:import_csv}
Il sistema deve consentire l'importazione di edifici da file in formato CSV. L'importazione deve avvenire in due fasi distinte: una fase di \emph{validazione a secco}, che non modifica il database e produce un'anteprima delle righe valide, delle righe scartate e delle relative motivazioni; e una fase di \emph{conferma}, che persiste esclusivamente le righe validate. L'header richiesto è \texttt{id,name,address,municipality,latitude,longitude,buildingType,surfaceM2}. Il sistema deve gestire i file con BOM UTF-8, con fine riga LF, CRLF o CR, con virgole all'interno di campi quotati, con separatore decimale punto o virgola, e deve ignorare le righe vuote. I campi obbligatori non possono essere vuoti; latitude e longitude devono essere numeriche; surfaceM2 deve essere numerica e strettamente positiva; gli identificativi già presenti nel catalogo devono essere rifiutati durante la validazione; gli identificativi ripetuti all'interno dello stesso file vengono scartati in fase di conferma. Il sistema non deve mai interpretare in modo silenzioso un valore malformato o mancante come zero: ogni riga scartata deve riportare la motivazione esplicita. Al termine della conferma il sistema deve riportare il numero di record inseriti, il numero di record scartati e il dettaglio degli scarti, segnalando eventuali incoerenze tra il conteggio degli scarti e il relativo dettaglio. Gli edifici importati devono essere marcati con provenienza sintetica.

\subsection{Requisiti funzionali di configurazione dello scenario}

\RF{Configurazione di uno scenario di ottimizzazione.} \label{rf:config_scenario}
Il sistema deve consentire all'operatore tecnico di costruire uno scenario di ottimizzazione specificando almeno: un budget complessivo disponibile in euro, un obiettivo da perseguire tra quelli supportati (risparmio energetico, risparmio economico, riduzione di CO\textsubscript{2}) e l'insieme delle opzioni di intervento candidate. Il sistema deve rendere sempre visibili all'operatore i parametri correnti dello scenario e un riepilogo con il costo degli interventi selezionati, il budget e il budget residuo, evidenziando quando il costo selezionato supera il budget. Il sistema può proporre valori di default demo (500.000 EUR e risparmio energetico) ma non deve applicare automaticamente un budget comunale reale né un obiettivo non visibile: i parametri correnti devono essere sempre mostrati e modificabili. Il sistema deve inoltre consentire di salvare il risultato corrente in uno dei due slot di confronto (A o B) al momento dell'esecuzione oppure dal dossier dei risultati.

\RF{Definizione dei vincoli dello scenario.} \label{rf:vincoli_scenario}
Il sistema deve consentire all'operatore di specificare sullo scenario i vincoli applicabili, tra cui:
\begin{itemize}[leftmargin=1.4em, itemsep=0.15em]
    \item coppie di opzioni mutuamente incompatibili;
    \item relazioni di prerequisito tra opzioni (un'opzione richiede che un'altra sia selezionata);
    \item un numero massimo di opzioni selezionabili all'interno di un sottoinsieme specificato;
    \item opzioni obbligatoriamente incluse oppure obbligatoriamente escluse dal portafoglio.
\end{itemize}
Il sistema deve mostrare in modo leggibile l'elenco dei vincoli applicati e deve impedire la costruzione di scenari logicamente incoerenti (ad esempio la stessa opzione contemporaneamente inclusa ed esclusa, un numero massimo non positivo oppure un riferimento a un'opzione non presente nel catalogo). Le incompatibilità sono simmetriche. Il sistema deve segnalare esplicitamente le situazioni in cui i vincoli rendono il problema inammissibile, senza eliminarli o ignorarli silenziosamente. Se i vincoli o i parametri dello scenario vengono modificati dopo un'esecuzione, il sistema deve rendere evidente che il risultato precedente non si riferisce più allo scenario corrente.

\RF{Selezione del sottoinsieme di candidati.} \label{rf:candidati}
Il sistema deve consentire di selezionare in modo puntuale quali opzioni di intervento entrano a far parte dell'insieme dei candidati dello scenario, a partire dall'intero catalogo. Il sistema deve fornire funzioni di selezione e deselezione rapida (seleziona tutti, deseleziona tutti), una ricerca per individuare rapidamente un intervento e una funzione di caricamento di un insieme dimostrativo di candidati. Il sistema deve mostrare in modo chiaro il numero di candidati correntemente selezionati, il costo complessivo della selezione e l'eventuale superamento del budget. Il sistema deve inoltre poter caricare un insieme dimostrativo di candidati con budget e obiettivo predefiniti, pensato esclusivamente per la demo e sempre modificabile dall'operatore.

\subsection{Requisiti funzionali di ottimizzazione e spiegazione}

\RF{Esecuzione dell'ottimizzazione.} \label{rf:ottimizzazione}
Il sistema deve, su richiesta esplicita dell'operatore, calcolare il portafoglio ottimo di opzioni di intervento rispetto all'obiettivo dello scenario, nel rispetto di tutti i vincoli specificati. Il sistema deve comunicare in modo esplicito lo stato del calcolo, distinguendo almeno: ottimo provato, ammissibile senza prova di ottimalità, inammissibile, indeterminato per timeout e modello non valido. In caso di soluzione disponibile, il sistema deve produrre la lista degli interventi selezionati e degli interventi esclusi, corredata degli indicatori aggregati dello scenario (investimento totale, budget residuo, risparmio energetico totale, risparmio economico totale, CO\textsubscript{2} evitata, numero di interventi selezionati) e del valore dell'obiettivo. Il calcolo deve essere deterministico: a parità di scenario, lo stesso portafoglio e gli stessi indicatori devono essere prodotti indipendentemente dall'ordine di presentazione dei candidati, come stabilito da \RNFref{rnf:determinismo}. Il sistema deve associare a ogni esecuzione un identificativo, un istante di esecuzione e un hash canonico del problema, e deve applicare criteri di spareggio deterministici in presenza di più soluzioni ottime equivalenti. Il sistema non deve presentare come definitivo un risultato prodotto in stato indeterminato.

\RF{Visualizzazione dei risultati dell'ottimizzazione.} \label{rf:risultati}
Il sistema deve presentare i risultati dell'ottimizzazione come un \emph{dossier} del portafoglio proposto, comprendente: una fascia di audit con identificativo dell'esecuzione, hash del problema, istante di esecuzione, stato del solver, indicazione dell'eventuale prova di ottimalità e stato del tie-break canonico; gli indicatori aggregati dello scenario con unità di misura esplicite; l'elenco dettagliato degli interventi selezionati con, per ciascuno, edificio, denominazione, investimento, risparmio energetico e motivazione; una rappresentazione dell'impiego del budget; l'allocazione dell'investimento per edificio; la possibilità di salvare il risultato negli slot A e B, di aprire il confronto e di stampare la relazione. Il dossier deve essere consultabile anche in un secondo momento all'interno della sessione, finché l'ultimo risultato è disponibile.

\RF{Spiegazione deterministica delle scelte.} \label{rf:spiegazioni}
Il sistema deve produrre, per ciascuna opzione candidata, una spiegazione deterministica e verificabile della sua selezione o esclusione, classificata in una delle seguenti categorie:
\begin{itemize}[leftmargin=1.4em, itemsep=0.15em]
    \item \emph{selezionato} per contributo diretto all'obiettivo;
    \item \emph{selezionato in quanto prerequisito} di altre opzioni selezionate;
    \item \emph{escluso per esclusione esplicita};
    \item \emph{escluso per incompatibilità} con un'opzione selezionata;
    \item \emph{escluso per prerequisito mancante};
    \item \emph{escluso per budget totale insufficiente};
    \item \emph{escluso per budget residuo insufficiente};
    \item \emph{escluso per limite numerico};
    \item \emph{escluso in quanto non ottimale} (causa specifica non deducibile localmente).
\end{itemize}
La spiegazione deve distinguere esplicitamente le cause dimostrabili localmente dalle cause non determinate, deve citare gli identificativi e i vincoli che la giustificano e non deve essere prodotta mediante generazione in linguaggio naturale non deterministica, come stabilito da \RNFref{rnf:interpretabilita}. Il sistema deve offrire una vista esplorabile delle motivazioni, con filtri per famiglia di causa, ricerca testuale e possibilità di mostrare tutte le opzioni.

\RF{Consultazione dello storico delle esecuzioni.} \label{rf:storico}
Il sistema deve registrare le esecuzioni di ottimizzazione in uno storico consultabile in sola lettura, mantenendo per ciascuna esecuzione i dati di audit disponibili, tra cui: identificativo dell'esecuzione, istante di esecuzione, scenario di riferimento, parametri dello scenario, stato del calcolo, indicatori aggregati e hash canonico del problema. Lo storico deve essere visibile nella dashboard come registro delle esecuzioni dello scenario e deve poter essere interrogato per scenario tramite le API del sistema; le esecuzioni non devono poter essere modificate o cancellate dall'interfaccia utente. Le analisi di sensibilità non devono essere registrate nello storico, perché non producono un risultato di audit: il sistema deve rendere esplicita questa distinzione.

\subsection{Requisiti funzionali di analisi comparativa}

\RF{Gestione degli scenari salvati.} \label{rf:scenari_salvati}
Il sistema deve consentire all'operatore di salvare il risultato corrente in due slot di confronto (A e B) e di rendere visibile in ogni momento se ciascuno slot è popolato o vuoto, con l'importo del risultato salvato. Il sistema deve consentire di richiamare uno scenario salvato per duplicarlo come nuovo scenario what-if, senza modificare l'originale, e di azzerare la sessione di lavoro (slot e ultimo risultato) previa conferma esplicita. Gli scenari salvati devono essere conservati per la durata della sessione del browser; il sistema deve segnalare in modo chiaro quando lo slot viene sovrascritto.

\RF{Confronto tra due scenari.} \label{rf:confronto}
Il sistema deve consentire di confrontare due scenari (A e B) e di mostrare, per ciascuno, i parametri principali e gli indicatori aggregati, mettendo in evidenza gli scostamenti calcolati come differenza $B - A$ su investimento, risparmio energetico, risparmio economico, CO\textsubscript{2} e numero di interventi selezionati. Il confronto deve includere una tabella delle metriche e due viste degli interventi aggiunti in B e rimossi rispetto ad A. Se i due scenari hanno obiettivi diversi, il sistema deve mostrare un avviso che chiarisce che gli scostamenti tra grandezze non omogenee vanno letti come differenza descrittiva e non come preferenza. Il sistema deve consentire di tornare al dossier dell'ultimo risultato e di stampare la vista di confronto.

\RF{Analisi di sensibilità sul budget.} \label{rf:sensibilita}
Il sistema deve consentire di eseguire un'analisi di sensibilità sul parametro budget, definendo un budget base, un obiettivo, il numero di punti della griglia (da 3 a 51) e i fattori minimo e massimo rispetto al budget base. L'analisi deve selezionare un insieme di opzioni candidate e rieseguire l'ottimizzazione in modo deterministico su ciascun punto della griglia, mantenendo invariati candidati, obiettivo e solver. Il sistema deve mostrare per ciascun punto il budget, lo stato del solver, l'investimento, il budget residuo, il valore dell'obiettivo e il numero di interventi selezionati; deve inoltre classificare ogni opzione in \emph{sempre selezionata}, \emph{mai selezionata} e \emph{oscillante} (selezionata solo per alcune soglie di budget), e rappresentare l'andamento dell'obiettivo e la mappa di selezione per budget. L'analisi deve essere dichiarata per quello che è: uno sweep deterministico su un solo parametro, non una simulazione Monte Carlo e non una previsione. Il sistema deve indicare esplicitamente che l'analisi varia solo il budget e che non applica gli altri vincoli dello scenario. Le esecuzioni dell'analisi non devono essere persistite nello storico di audit.

\subsection{Requisiti funzionali di assistenza AI (opzionali)}

\RF{Proposta assistita di configurazione dello scenario.} \label{rf:ai_proposta}
Il sistema deve offrire, in modo opzionale, la possibilità di ottenere una proposta di configurazione dello scenario a partire da una richiesta in linguaggio naturale dell'operatore (ad esempio ``\emph{vorrei concentrarmi sulla riduzione di CO\textsubscript{2} con 500 mila euro}''). La proposta deve riguardare esclusivamente i parametri di configurazione (budget, obiettivo, insieme dei candidati) e deve essere restituita in forma strutturata, corredata da una motivazione, dall'indicazione dell'origine (motore a regole locale, modello linguistico esterno oppure fallback dal modello esterno al motore locale), da eventuali informazioni mancanti e da eventuali avvertenze. La proposta non deve introdurre identificativi di opzioni non presenti nel catalogo e non deve calcolare risultati. L'applicazione deve essere pienamente funzionante in assenza di configurazione di un provider linguistico esterno.

\RF{Approvazione umana della proposta AI.} \label{rf:ai_approvazione}
Il sistema deve richiedere sempre un atto esplicito di approvazione da parte dell'operatore prima di applicare una proposta AI allo scenario corrente. Il sistema deve mostrare in modo evidente che la proposta non è stata applicata e deve consentire all'operatore di accettarla, modificarla oppure rifiutarla; in nessun caso la proposta AI deve essere applicata automaticamente. In caso di accettazione, il sistema deve aggiornare budget, obiettivo e insieme dei candidati, mostrare un avviso che ricorda di verificare i parametri e lasciare all'operatore il compito di avviare l'ottimizzazione. Le opzioni che la proposta suggerisce di escludere devono essere evidenziate senza essere deselezionate automaticamente. In caso di rifiuto, lo scenario corrente deve restare invariato.

\subsection{Requisiti funzionali di presentazione e trasparenza}

\RF{Cambio del registro di visualizzazione e stampa della relazione.} \label{rf:tema}
Il sistema deve offrire due registri visivi dello stesso strumento: \emph{Relazione}, pensato per la lettura diurna, la proiezione e la stampa, e \emph{Sala operativa}, pensato per sessioni prolungate di analisi su schermo. I due registri devono condividere il medesimo layout e la medesima informazione, variando esclusivamente il materiale cromatico e tipografico. Il sistema deve consentire all'utente di passare da un registro all'altro in qualunque momento, deve conservare la scelta tra una sessione e la successiva e deve utilizzare la preferenza del sistema operativo come valore di default alla prima apertura. Il sistema deve consentire di stampare la relazione del risultato e la vista di confronto; in stampa il sistema deve produrre un documento a fondo chiaro con inchiostro scuro, nascondendo gli elementi di navigazione e i controlli non pertinenti. Il sistema deve inoltre offrire una palette di comandi rapidi (Ctrl/Cmd+K) per raggiungere le sezioni e attivare le azioni principali, tra cui il cambio registro e la stampa.

\RF{Trasparenza sulla provenienza dei dati.} \label{rf:provenienza}
Il sistema deve rendere sempre riconoscibile la natura dei dati utilizzati, come stabilito da \RNFref{rnf:trasparenza}. In particolare: deve mostrare, quando lo spazio lo consente, il timbro che dichiara il dataset demo e sintetico e che nessun dato reale del Comune di Trento è incluso; deve mostrare la provenienza per ciascun edificio e per ciascuna opzione di intervento; deve riportare nel piè di pagina la fonte del dataset, il modello di ottimizzazione adottato e l'istante dell'ultima esecuzione; deve etichettare il payback come stima demo; non deve presentare valori sintetici come misure ufficiali. Nei viewport più piccoli e in stampa il timbro può essere compattato, ma il piè di pagina di provenienza deve restare disponibile.
